\documentclass[10pt, letterpaper]{article}

% Packages:
\usepackage[
    ignoreheadfoot, % set margins without considering header and footer
    top=2 cm, % seperation between body and page edge from the top
    bottom=2 cm, % seperation between body and page edge from the bottom
    left=2 cm, % seperation between body and page edge from the left
    right=2 cm, % seperation between body and page edge from the right
    footskip=1.0 cm, % seperation between body and footer
    % showframe % for debugging 
]{geometry} % for adjusting page geometry
\usepackage{titlesec} % for customizing section titles
\usepackage{tabularx} % for making tables with fixed width columns
\usepackage{array} % tabularx requires this
\usepackage[dvipsnames]{xcolor} % for coloring text
\definecolor{primaryColor}{RGB}{0, 0, 0} % define primary color
\usepackage{enumitem} % for customizing lists
\usepackage{fontawesome5} % for using icons
\usepackage{amsmath} % for math
\usepackage[
    pdftitle={John Doe's CV},
    pdfauthor={John Doe},
    pdfcreator={LaTeX with RenderCV},
    colorlinks=true,
    urlcolor=primaryColor
]{hyperref} % for links, metadata and bookmarks
\usepackage[pscoord]{eso-pic} % for floating text on the page
\usepackage{calc} % for calculating lengths
\usepackage{bookmark} % for bookmarks
\usepackage{lastpage} % for getting the total number of pages
\usepackage{changepage} % for one column entries (adjustwidth environment)
\usepackage{paracol} % for two and three column entries
\usepackage{ifthen} % for conditional statements
\usepackage{needspace} % for avoiding page brake right after the section title
\usepackage{iftex} % check if engine is pdflatex, xetex or luatex

% Ensure that generate pdf is machine readable/ATS parsable:
\ifPDFTeX
    \input{glyphtounicode}
    \pdfgentounicode=1
    \usepackage[T1]{fontenc}
    \usepackage[utf8]{inputenc}
    \usepackage{lmodern}
\fi

\usepackage{charter}

% Some settings:
\raggedright
\AtBeginEnvironment{adjustwidth}{\partopsep0pt} % remove space before adjustwidth environment
\pagestyle{empty} % no header or footer
\setcounter{secnumdepth}{0} % no section numbering
\setlength{\parindent}{0pt} % no indentation
\setlength{\topskip}{0pt} % no top skip
\setlength{\columnsep}{0.15cm} % set column seperation
\pagenumbering{gobble} % no page numbering

\titleformat{\section}{\needspace{4\baselineskip}\bfseries\large}{}{0pt}{}[\vspace{1pt}\titlerule]

\titlespacing{\section}{
    % left space:
    -1pt
}{
    % top space:
    0.3 cm
}{
    % bottom space:
    0.2 cm
} % section title spacing

\renewcommand\labelitemi{$\vcenter{\hbox{\small$\bullet$}}$} % custom bullet points
\newenvironment{highlights}{
    \begin{itemize}[
        topsep=0.10 cm,
        parsep=0.10 cm,
        partopsep=0pt,
        itemsep=0pt,
        leftmargin=0 cm + 10pt
    ]
}{
    \end{itemize}
} % new environment for highlights


\newenvironment{highlightsforbulletentries}{
    \begin{itemize}[
        topsep=0.10 cm,
        parsep=0.10 cm,
        partopsep=0pt,
        itemsep=0pt,
        leftmargin=10pt
    ]
}{
    \end{itemize}
} % new environment for highlights for bullet entries

\newenvironment{onecolentry}{
    \begin{adjustwidth}{
        0 cm + 0.00001 cm
    }{
        0 cm + 0.00001 cm
    }
}{
    \end{adjustwidth}
} % new environment for one column entries

\newenvironment{twocolentry}[2][]{
    \onecolentry
    \def\secondColumn{#2}
    \setcolumnwidth{\fill, 4.5 cm}
    \begin{paracol}{2}
}{
    \switchcolumn \raggedleft \secondColumn
    \end{paracol}
    \endonecolentry
} % new environment for two column entries

\newenvironment{threecolentry}[3][]{
    \onecolentry
    \def\thirdColumn{#3}
    \setcolumnwidth{, \fill, 4.5 cm}
    \begin{paracol}{3}
    {\raggedright #2} \switchcolumn
}{
    \switchcolumn \raggedleft \thirdColumn
    \end{paracol}
    \endonecolentry
} % new environment for three column entries

\newenvironment{header}{
    \setlength{\topsep}{0pt}\par\kern\topsep\centering\linespread{1.5}
}{
    \par\kern\topsep
} % new environment for the header

\newcommand{\placelastupdatedtext}{% \placetextbox{<horizontal pos>}{<vertical pos>}{<stuff>}
  \AddToShipoutPictureFG*{% Add <stuff> to current page foreground
    \put(
        \LenToUnit{\paperwidth-2 cm-0 cm+0.05cm},
        \LenToUnit{\paperheight-1.0 cm}
    ){\vtop{{\null}\makebox[0pt][c]{
        \small\color{gray}\textit{Last updated in September 2024}\hspace{\widthof{Last updated in September 2024}}
    }}}%
  }%
}%

% save the original href command in a new command:
\let\hrefWithoutArrow\href

% new command for external links:


\begin{document}
    \newcommand{\AND}{\unskip
        \cleaders\copy\ANDbox\hskip\wd\ANDbox
        \ignorespaces
    }
    \newsavebox\ANDbox
    \sbox\ANDbox{$|$}

    \begin{header}
        \fontsize{25 pt}{25 pt}\selectfont Trunov Vladimir

        \vspace{5 pt}

        \normalsize
        \kern 5.0 pt%
        \mbox{\href{https://telegram.me/vvtrunov}{\underline{Telegram}}}
        \kern 5.0 pt%
        \AND
        \kern 5.0 pt%
        \mbox{\hrefWithoutArrow{mailto:trunov.vv@phystech.edu}{\underline{trunov.vv@phystech.edu}}}%
        \kern 5.0 pt%
        \AND%
        \kern 5.0 pt%
        \mbox{\hrefWithoutArrow{tel:+7 (926) 692-03-22}{+7 (926) 692-03-22}}%
        \kern 5.0 pt%

    \end{header}

    \vspace{5 pt - 0.3 cm}

    \section{Education}
        \begin{twocolentry}{
            Sept 2021 – Present
        }
            \textbf{Moscow institute of physics and technology}, Physics and Computer Science\end{twocolentry}

        \vspace{0.10 cm}
        \begin{onecolentry}
            \begin{highlights}
                \item \textbf{Bachelor's degree thesis:} Development of a domain-specific language for
mapping nodes of the abstract syntactic
tree of the SystemVerilog hardware description language
            \end{highlights}
        \end{onecolentry}


    
    \section{Experience}
    \begin{twocolentry}{
            July 2022 – Dec 2023
        }
            \textbf{C++ Software Engineer}, Compilers Department, ISP RAS -- Moscow, Russia \end{twocolentry}
    
        \vspace{0.10 cm}

        \begin{itemize}[label={}]

        \item \begin{twocolentry} {
            July 2022 - June 2023
        }
            \textbf{{Clang-Tidy/Clang Static Analyzer}}
        \end{twocolentry}

        Developed C/C++ source code vulnerability analysis tools based on Clang-Tidy and Clang Static
Analyzer. Implemented a set of static checks according to AUTOSAR C++ 14 and Misra-C standards.

        \item \begin{twocolentry} {
            July 2023 - Dec 2023
        }
            \textbf{{LLVM Backend}}
        \end{twocolentry}
I wrote patches to the Clang compiler for the LLVM Backend extension to support
code generation for a new processor based on the PowerPC architecture. I also added support from third-party
libraries and tools such as Address Sanitizer and OpenMP.

        \end{itemize}

        \vspace{0.4 cm}
\begin{twocolentry}{
            Jan 2024 – March 2025
        }
            \textbf{Software Engineer}, Yandex -- Moscow, Russia\end{twocolentry}
            \vspace{0.10 cm}
I worked in the Yandex.Market Search and data preparation infrastructure department. My main language was C++, sometimes I had to use Python. I was engaged in the following
tasks:
        \vspace{0.10 cm}
        \begin{onecolentry}
            \begin{highlights}
                \item Optimized the storage scheme in the database for fast search and low access timings.
                \item I rewrote the legacy big-data service, increasing its performance by 2-3 times.
                \item I have written a new microservice that increases the fault tolerance of the system.
            \end{highlights}
        \end{onecolentry}

        \vspace{0.4 cm}
        \begin{twocolentry}{
            March 2025 – Dec 2025
        }
            \textbf{Research Software Engineer}, Kaspersky -- Moscow, Russia\end{twocolentry}
            \vspace{0.10 cm}
I worked in the department of testing and performance analysis. My responsibilities included:
        \vspace{0.10 cm}
        \begin{onecolentry}
            \begin{highlights}
                \item Profiling of Linux products using Linux Perf and various eBPF scripts.
                \item Analysis and optimization of product architecture.
                \item Development of internal infrastructure for automation of testing and profiling.
            \end{highlights}
        \end{onecolentry}

        \vspace{0.4 cm}
        \begin{twocolentry}{
            January 2026 – Present
        }
            \textbf{Software Engineer}, Arenadata \end{twocolentry}
            \vspace{0.10 cm}
    My team developing new open source storage engine OrioleDB for PostgreSQL. As a developer my responsibilities include:
    \begin{highlights}
        \item Implementing functionalities from the vanilla PostgreSQL engine, I implemented several DDL features and fixed some bugs.
        \item Profiling using Linux-based profilers. I have been involved in researching performance issues, and in this particular case, we achieved a 40x decrease in disk usage.
    \end{highlights}
            \vspace{0.2 cm}


    \section{Technologies}
        \begin{onecolentry}
            \textbf{Languages:} C++, C, Python \end{onecolentry}
        \begin{onecolentry}
            \textbf{Tools:} Linux, clang/gcc, make, CMake, gdb/lldb, git, Linux Perf, Valgrind, eBPF
\end{onecolentry}
\end{document}
